\chapter{Planificación temporal}

El desarrollo del Trabajo de Fin de Grado se ha llevado a cabo a lo largo de un periodo aproximado de quince meses, compatibilizando su realización con una actividad laboral a tiempo completo. Esta circunstancia ha condicionado la dedicación semanal y ha dado lugar a una planificación flexible y progresiva, con fases de menor y mayor intensidad según la carga de trabajo y la madurez del proyecto.

El proyecto ha seguido una metodología iterativa e incremental, permitiendo redefinir el alcance inicial y adaptar los objetivos conforme se consolidaban las decisiones técnicas.

\section{Fases del proyecto}

Con el objetivo de organizar el desarrollo del Trabajo de Fin de Grado de manera estructurada y coherente, el proyecto se ha dividido en una serie de fases claramente diferenciadas. Cada una de estas fases responde a un conjunto de objetivos específicos y refleja la evolución progresiva del trabajo, desde la definición inicial del alcance hasta la validación final del sistema y la elaboración de la documentación.

\subsection{Fase 1: Exploración inicial y definición de alcance (meses 1--3)}

Durante esta fase se realizó una exploración inicial de posibles enfoques para el TFG, evaluando distintas ideas de proyecto y tecnologías. Se analizaron los requisitos generales del trabajo académico y se llevó a cabo una primera toma de contacto con el desarrollo de videojuegos en Unity.

Como resultado de esta fase se decidió redefinir el enfoque del TFG hacia el desarrollo de un videojuego del subgénero \textit{rogue-like}, al considerarse un dominio adecuado para aplicar conocimientos de ingeniería del software.

\subsection{Fase 2: Diseño conceptual y arquitectura base (meses 4--6)}

En esta fase se definieron los objetivos generales y específicos del proyecto, así como la arquitectura inicial del sistema. Se diseñaron los principales subsistemas del juego (gestión de estados, flujo de salas, generación procedural y combate) y se establecieron las bases del proyecto en Unity.

Se tomaron decisiones clave relativas al uso de \textit{ScriptableObjects}, patrones de diseño y separación entre lógica y datos.

\subsection{Fase 3: Implementación progresiva del núcleo del juego (meses 7--11)}

Esta fase constituyó el núcleo del trabajo práctico. Se implementaron de forma incremental los sistemas fundamentales del videojuego:

\begin{itemize}
    \item Generación procedural de salas
    \item Gestión del flujo de juego
    \item Sistema de combate y armas
    \item Inteligencia artificial de enemigos
    \item Sistema de jefes y condición de victoria.
\end{itemize}

La implementación se realizó de manera iterativa, validando la funcionalidad del juego en cada incremento y ajustando decisiones de diseño cuando fue necesario.

\subsection{Fase 4: Herramientas de editor y mejora de productividad (meses 10--13)}

De forma paralela a la implementación del juego, se desarrollaron utilidades internas de apoyo al desarrollo dentro del editor de Unity con el objetivo de automatizar tareas repetitivas y facilitar la creación y validación de contenido.

Estas herramientas permitieron acelerar el desarrollo y reducir errores, especialmente en la configuración de salas, enemigos y armas.

\subsection{Fase 5: Refinamiento, pruebas y documentación (meses 14--15)}

En la fase final se realizó un proceso de pruebas funcionales del sistema completo, verificando la estabilidad del flujo de juego y corrigiendo errores detectados. Se llevaron a cabo ajustes de balance y mejoras en la interfaz de usuario.

Finalmente, se procedió a la redacción de la memoria académica, documentando el proceso de desarrollo, las decisiones técnicas adoptadas y las limitaciones encontradas.
